\documentclass[11pt]{article}
\usepackage[margin=2.5cm]{geometry}
\usepackage[utf8]{inputenc}
%\usepackage[spanish]{babel}
\usepackage{amsmath,amssymb,booktabs}
\usepackage{hyperref}
\newcommand{\B}{\mathbf{B}}
\newcommand{\Lm}{\mathbf{L}}
\newcommand{\Dm}{\mathbf{D}}
\newcommand{\X}{\mathbf{X}}
\title{Localización por asignación de observaciones:\\ qué hace el método, qué falló y qué lo corrige}
\author{Nota de trabajo}
\date{14 de septiembre de 2026}

\begin{document}\maketitle

\section{Qué hace el método (esto no cambia)}

Hay $n$ componentes del estado y $p$ observaciones. En cada ciclo:

\begin{enumerate}
\item Cada componente $i$ elige \emph{una} observación $a(i)$ entre las que tiene cerca.
\item Su radio de localización es la distancia a esa observación:
  $r_i = \max\big(1,\ \mathrm{dist}(i, a(i))\big)$.
  El piso $1$ evita que un componente observado en su propio sitio quede con radio $0$, que
  dejaría su fila de la precisión diagonal (medido: eso solo cuesta un 25\,\% de error).
\item Con esos radios se arma la precisión de Cholesky modificado: el componente $i$ se regresa
  sobre sus predecesores $P(i)$, los vecinos con índice menor a distancia $\le r_i$,
  \begin{equation}
    \boldsymbol\beta_i = \big(\X_P^\top\X_P + a_i I\big)^{-1}\X_P^\top\,\Delta\mathbf{x}_i,
    \qquad a_i = \alpha\,\frac{\operatorname{tr}(\X_P^\top\X_P)}{|P(i)|},
  \end{equation}
  y $\Lm_{ii}=1$, $\Lm_{i,P(i)} = -\boldsymbol\beta_i^\top$, $\Dm_{ii} = 1/\mathrm{var(residuo)}$,
  $\widehat\B^{-1} = \Lm^\top\Dm\Lm$.
\item El análisis es el global de siempre, con \emph{todas} las observaciones:
  \begin{equation}
    \big(\widehat\B^{-1} + \mathbf{H}^\top\mathbf{R}^{-1}\mathbf{H}\big)\,\delta\mathbf{x}
      = \mathbf{H}^\top\mathbf{R}^{-1}(\mathbf{y}-\mathbf{H}\mathbf{x}^b).
  \end{equation}
\end{enumerate}

La asignación decide \emph{radios}, no qué observaciones corrigen a quién. Toda la información
entra en el paso 4. Lo único que cambia entre versiones del método es el paso 1: la regla con
la que $i$ elige $a(i)$.

\section{Qué falló: el costo $J$}

La regla original elegía $a(i)$ minimizando el costo variacional de la actualización escalar
con la observación $j$, con $\sigma_{ij}$ la covarianza del ensamble entre $i$ y el sitio $j$,
$\sigma_{jj}$ la varianza allí, $r$ el error de observación y $\delta_j = y_j - x^b_j$ la
innovación:
\begin{equation}
  J_i(j) = \frac{(x^a_i - x^b_i)^2}{\sigma_{ii}} + \frac{(y_j - x^a_j)^2}{r}
  \;=\; \frac{\delta_j^2}{(\sigma_{jj}+r)^2}\left[\frac{\sigma_{ij}^2}{\sigma_{ii}} + r\right].
  \label{eq:J}
\end{equation}
La forma factorizada muestra el problema. El factor $\delta_j^2/(\sigma_{jj}+r)^2$ depende solo del
sitio $j$, no de $i$. Lo único que depende de $i$ es el corchete, y se minimiza cuando
$\sigma_{ij}^2$ es \textbf{lo más pequeño posible}: $J$ prefiere, entre las candidatas, la
observación \emph{menos} correlacionada con $i$, porque es la que menos lo mueve. Ningún término
premia que $x_i$ mejore; no puede, porque $x_i$ no está observado.

Medido en el QG: la correlación mediana de la observación elegida era $0{,}15$ cuando la
máxima disponible era $0{,}48$; $J$ elegía la de menor correlación el 44\,\% de las veces. El
``79\,\% que no toma la más cercana'' era esto: elegían por inercia, no por información.

En el análisis global (paso 4) el daño quedaba escondido, porque las demás observaciones
compensan un radio mal elegido. Se hizo visible al restringir la actualización a la observación
asignada: ahí el análisis con $J$ dejaba el error igual que el fondo.

\section{El primer intento de corrección, y por qué tampoco basta}

Lo natural es elegir la observación que más reduce la varianza del error de análisis en $i$:
\begin{equation}
  a(i) = \arg\max_j \frac{\sigma_{ij}^2}{\sigma_{jj}+r}.
  \label{eq:var}
\end{equation}
Es el criterio correcto en el modelo lineal-gaussiano y no usa la innovación. Corrige la
dirección: ahora sí elige las correlacionadas. Pero $\sigma_{ij}$ es la covarianza \emph{cruda}
del ensamble, y con $N=40$ miembros dos puntos no correlacionados muestran una correlación
espuria del orden de $1/\sqrt{N}\approx 0{,}16$. Entre muchas candidatas, alguna lejana siempre
tiene ruido mayor que la señal de la vecina real: los radios salen grandes y ruidosos
(medio $2{,}9$, máximo $8$) y el análisis pierde contra el radio fijo ($0{,}381$ contra
$0{,}320$). Es el mismo problema por el que la localización existe.

\section{La corrección que funciona: leer la elección de $\widehat\B^{-1}$}

En lugar de la covarianza cruda, la dependencia \emph{condicional} que el propio Cholesky
estima. En cada ciclo, antes del paso 1:

\begin{enumerate}
\item[0.] Armar una precisión de tanteo $\widehat\B^{-1}_{\rho}=\Lm^\top\Dm\Lm$ con un radio
  ancho $\rho$ igual para todos (sin invertir nada; es rala y sus entradas están disponibles).
\end{enumerate}
La entrada $(\widehat\B^{-1}_\rho)_{ij}$ mide cuánto depende $i$ del sitio $j$ \emph{dado todo lo
demás}. Normalizada es la correlación parcial,
\begin{equation}
  \varrho_{ij} = -\frac{(\widehat\B^{-1}_\rho)_{ij}}
                       {\sqrt{(\widehat\B^{-1}_\rho)_{ii}\,(\widehat\B^{-1}_\rho)_{jj}}},
  \qquad
  a(i) = \arg\max_{j \text{ observado}} |\varrho_{ij}|.
  \label{eq:partial}
\end{equation}
Después, pasos 1--4 como siempre: $r_i = \max(1,\mathrm{dist}(i,a(i)))$, se rearma
$\widehat\B^{-1}$ con esos radios, análisis global.

Por qué esto sí: la regresión con ridge ya descontó lo que los vecinos explican, así que una
correlación cruda espuria entre $i$ y un punto lejano no sobrevive como dependencia condicional.
El criterio y el estimador de $\B^{-1}$ son el mismo objeto: se elige el radio de cada fila con
la información que la propia fila produce.

\paragraph{El radio de tanteo $\rho$ no es libre.} Tiene que ser uno que la regresión aguante
con $N$ miembros. Con $\rho = 8$ cada fila tiene $\sim 150$ predecesores para $40$ miembros, los
coeficientes son ruido y todo se cae (radio medio $4{,}8$, error $0{,}68$). Con $\rho=4$ son
$\sim 40$ predecesores y funciona. Es el único parámetro nuevo, y va atado a $N$.

\section{Los números (QG, $N=40$, retícula de paso 2, 10 ciclos, una semilla)}

Error RMS del análisis en $q$ sobre el interior, unidades normalizadas, media de los ciclos 4--9.
Todos los brazos con el mismo análisis global EnKF-MC.

\begin{center}\begin{tabular}{@{}lccr@{}}\toprule
regla de radio & radio medio & radio máx & RMSE \\ \midrule
fijo $r=1$ & 1 & 1 & 0,322 \\
fijo $r=2$ (mejor uniforme) & 2 & 2 & 0,320 \\
fijo $r=4$ & 4 & 4 & 0,468 \\ \midrule
asignación con $J$ \eqref{eq:J} & 1,49 & 2 & 0,331 \\
asignación por varianza \eqref{eq:var}, $\kappa=4$, sin piso & 0,94 & 4 & 0,446 \\
asignación por varianza, candidatas a $\le 8$, piso 1 & 2,90 & 8 & 0,381 \\ \midrule
\textbf{asignación por correlación parcial \eqref{eq:partial}, $\rho=4$} & 2,01 & 7 & \textbf{0,3135} \\
\quad idem, $\rho=2$ & 1,23 & 4 & 0,321 \\
\quad idem, $\rho=8$ & 4,77 & 14 & 0,681 \\ \bottomrule
\end{tabular}\end{center}

Con $\rho=4$ la asignación queda por delante del mejor radio fijo ($0{,}3135$ contra $0{,}320$,
un 2\,\%). Es una semilla y diez ciclos: una señal, no un resultado. Es la primera vez que una
regla de asignación queda por delante del barrido uniforme.

\section{Qué sigue}

El benchmark ya montado en el repositorio (dos filtros, dos redes, $N\in\{40,80,120\}$, tres
semillas, 60 ciclos), con la regla \eqref{eq:partial} en lugar de $J$, y el radio de tanteo
$\rho\in\{3,4,5\}$ como eje. Hay que corregir \texttt{FINDINGS.md} y el paper: el hallazgo del
79\,\% es correcto como dato y falso como interpretación.

\end{document}
